\section{Cosmos and Mind}

``Philosophers'' of naturalism, in a sort of Heideggerian idle chatter, create a consensus group illusion through their articles, books, journals, blogs, and so on. They have their own ``sages'', call each other ``bright'' as a mark of mutual recognition, and rant against the ``sky god'' or ``flying spaghetti monsters''. So when one of their own apostatizes, as apparently \textbf{Thomas Nagel} has in his recent book \textit{Mind and Cosmos}\footnote{\url{https://www.amazon.com/Mind-Cosmos-Materialist-Neo-Darwinian-Conception/dp/0199919755/ref=sr\_1\_1\%20?tag=firstthings-20-20}}, there is a flurry of activity to restore the \emph{status quo ante}. On the one hand, the consensus illusion\footnote{\url{http://whyevolutionistrue.wordpress.com/2012/10/13/philosopher-thomas-nagel-goes-the-way-of-alvin-plantinga-disses-evolution/}} needs to be re-fortified. On the other, religious philosophers such as Edward Feser\footnote{\url{https://edwardfeser.blogspot.com/2012/10/nagel-and-his-critics-part-i.html}} or the Maverick Philosopher\footnote{\url{http://maverickphilosopher.typepad.com/maverick\_philosopher/nagel-thomas/}} are excited to spot a lost sheep, who, even if it hasn't quite returned to the fold, at least it has come out of hiding.

Mr. Nagel, considered one of the most prominent philosophers of our time, thinks he has ``discovered'' some important ideas. Of course, his ``discoveries'' have been known since the hoary times of the \textbf{Primordial Tradition}. Every Medieval monk knew what Mr. Nagel had not known. So much for ``progress'' from the ``dark ages''. It may be useful to quickly reflect on the five principles recently ``discovered'' by Mr. Nagel, as outlined by the Maverick Philosopher.

First of all, Mr. Nagel acknowledges that reductive materialism is untenable and consciousness---let us call it what it is, \emph{viz.}, Spirit---is real. As Gornahoor has repeated often, form and matter, Purusha and Prakriti, are necessary conditions of existence, so there is absolutely no surprise at these two discoveries. For the Primordial Tradition, matter (Prakriti) without form (Purusha) is pure potentiality, i.e., nothing in particular. In quantum physical terms, prime matter is a Schrödinger equation in an indeterminate state. A conscious agent is required to ``collapse'' the equation, actualizing thereby one possibility rather than another.

From a purely philosophical point of view, there may be some frustration in this, given the interminable debates about whether forms are real or not. In the Primordial Tradition, however, the forms are intuited, not rationally proven. Hence, one is not properly a metaphysician until he has developed that intellectual intuition that surpasses reason.

Mr. Nagel's more interesting claim is, ``Nature is intelligible. Its intelligibility is inherent in it.'' Readers of Gornahoor will immediately understand what he is getting at. Taking ``nature'' as referring to ``being'' (a term a naturalist eschews) and intelligibility as ``Logos'', we see that he has inchoately formed the notion of the identity of Being and Logos. Again, this is Tradition and, \emph{a fortiori}, is one of the principle creedal dogmas of the dominant spiritual tradition in the West for the past 2000 years.

Hence, \emph{pace} the Kantians and post-modernists, the cosmos itself is intelligible; it is therefore not the case than man imposes his own intelligence on some primal chaos. We can praise Mr. Nagel, I suppose, for rediscovering sliced bread, and we hope it doesn't destroy his reputation. Gornahoor, on the other hand, is a remote outpost on the Internet and has no reputation. Rightly so, since we have no credentials and make no claim to originality. We just hand down what has been handed down to us. Others may grab the baton or fumble it, that is not our concern. We just hope that readers avoid being distracted by alleged novelties.


\flrightit{Posted on 2012-10-24 by Cologero}

